% ============================================================
% Chapter IV — Results and Discussion
% ============================================================
\section{Results and Discussion}

\subsection{Box's M Test: Homogeneity of Covariance Matrices}

The Box's M test was conducted to determine whether the variance--covariance
matrices of the continuous predictors (Grid Position, Tire Age, Fuel Mass
Proxy, and Degradation $\times$ Fuel) were equal across the three finishing
tiers.

\begin{table}[htbp]
  \centering
  \caption{Box's M Test Results}
  \label{tab:boxs-m}
  \begin{tabular}{lc}
    \toprule
    \textbf{Statistic} & \textbf{Value} \\
    \midrule
    Chi-square (approx.) & 160.419 \\
    Degrees of freedom   & 20 \\
    $p$-value            & $< 2.2 \times 10^{-16}$ \\
    \bottomrule
  \end{tabular}
\end{table}

The test yielded $\chi^{2}(20) = 160.419$, $p < 2.2 \times 10^{-16}$,
formally rejecting the null hypothesis of equal covariance matrices across
tiers. This result has direct methodological implications: Multiple
Discriminant Analysis (MDA) requires homogeneous covariance matrices for valid
classification, and the observed violation disqualifies it as a candidate
method. \textbf{Multinomial logistic regression} was selected as the primary
inferential model because it does not require this assumption.

\subsection{VIF Diagnostics: Multicollinearity Check}

Before fitting the multinomial model, variance inflation factors were computed
for all predictors to ensure that the engineered features---particularly the
interaction term---did not introduce problematic multicollinearity.

\begin{table}[htbp]
  \centering
  \caption{Variance Inflation Factor (VIF) Diagnostics}
  \label{tab:vif}
  \begin{tabular}{lc}
    \toprule
    \textbf{Variable} & \textbf{VIF} \\
    \midrule
    GridPosition       & 1.007 \\
    Tire\_Age          & 3.220 \\
    Fuel\_Mass\_Proxy  & 2.315 \\
    Degradation\_x\_Fuel & 1.725 \\
    Tire\_MEDIUM       & 1.181 \\
    Tire\_SOFT         & 1.156 \\
    \bottomrule
  \end{tabular}
\end{table}

All VIF values fall well below the standard threshold of 10, with the highest
being 3.22 for Tire\_Age. The mean-centering applied to the interaction term
effectively prevented structural multicollinearity---the Degradation $\times$
Fuel VIF of 1.725 confirms that the interaction product and its components
remain statistically independent. Notably, the near-perfect VIF of 1.007 for
GridPosition confirms that starting position and tire degradation are
mathematically unrelated, which is physically expected since grid position is
determined in qualifying while tire age accumulates during the race.

\subsection{Multinomial Logistic Regression (MLE)}

The multinomial logistic regression model was fitted with ``No Points'' as the
baseline category. All coefficients represent log-odds relative to this
reference tier.

\subsubsection{Model Fit Statistics}

\begin{table}[htbp]
  \centering
  \caption{Multinomial Logistic Regression Model Fit}
  \label{tab:model-fit}
  \begin{tabular}{lc}
    \toprule
    \textbf{Statistic} & \textbf{Value} \\
    \midrule
    Residual Deviance & 3821.657 \\
    AIC               & 3849.657 \\
    Observations      & 2,880 \\
    \bottomrule
  \end{tabular}
\end{table}

\subsubsection{Coefficient Estimates (Log-Odds)}

\begin{table}[htbp]
  \centering
  \caption{MLE Coefficient Estimates (Log-Odds)}
  \label{tab:coefficients}
  \small
  \begin{tabular}{lrr}
    \toprule
    \textbf{Predictor} & \textbf{Podium} & \textbf{Point Scorer} \\
    \midrule
    Intercept           &  4.8565  &  2.9230 \\
    GridPosition        & $-0.7978$ & $-0.2271$ \\
    Tire\_Age           &  0.0291  & $-0.0036$ \\
    Fuel\_Mass\_Proxy   & $-0.0149$ & $-0.0084$ \\
    Degradation\_x\_Fuel &  0.0008  &  0.0002 \\
    Tire\_MEDIUM        &  0.2267  & $-0.4601$ \\
    Tire\_SOFT          &  0.2424  & $-0.2538$ \\
    \bottomrule
  \end{tabular}
\end{table}

\subsubsection{Standard Errors}

\begin{table}[htbp]
  \centering
  \caption{Standard Errors of Coefficient Estimates}
  \label{tab:std-errors}
  \small
  \begin{tabular}{lrr}
    \toprule
    \textbf{Predictor} & \textbf{Podium} & \textbf{Point Scorer} \\
    \midrule
    Intercept           & 0.4206  & 0.2511 \\
    GridPosition        & 0.0347  & 0.0107 \\
    Tire\_Age           & 0.0136  & 0.0075 \\
    Fuel\_Mass\_Proxy   & 0.0045  & 0.0023 \\
    Degradation\_x\_Fuel & 0.0003  & 0.0002 \\
    Tire\_MEDIUM        & 0.2628  & 0.1376 \\
    Tire\_SOFT          & 0.2740  & 0.1317 \\
    \bottomrule
  \end{tabular}
\end{table}

\subsubsection{Exact $p$-Values}

\begin{table}[htbp]
  \centering
  \caption{$p$-Values for Coefficient Significance}
  \label{tab:pvalues}
  \small
  \begin{tabular}{lrr}
    \toprule
    \textbf{Predictor} & \textbf{Podium} & \textbf{Point Scorer} \\
    \midrule
    Intercept           & 0.0000 & 0.0000 \\
    GridPosition        & 0.0000 & 0.0000 \\
    Tire\_Age           & 0.0328 & 0.6337 \\
    Fuel\_Mass\_Proxy   & 0.0009 & 0.0003 \\
    Degradation\_x\_Fuel & 0.0095 & 0.2174 \\
    Tire\_MEDIUM        & 0.3883 & 0.0008 \\
    Tire\_SOFT          & 0.3764 & 0.0540 \\
    \bottomrule
  \end{tabular}
\end{table}

\subsubsection{Interpretation of Key Findings}

Two results stand out from the significance testing:

\begin{enumerate}
  \item \textbf{GridPosition ($p = 0.000$ for both tiers).} Starting position
    is overwhelmingly significant, confirming its role as a baseline control.
    The coefficient of $-0.80$ for the Podium tier means that each position
    further back on the grid reduces the log-odds of a podium finish by
    approximately 0.80 units---a substantial effect that reflects the
    difficulty of overtaking in modern F1.

  \item \textbf{Degradation $\times$ Fuel ($p = 0.0095$ for Podium).} With
    grid position controlling for raw car pace, the interaction term achieves
    significance at $\alpha < 0.01$ for the Podium tier. The positive
    coefficient indicates that when tire age and fuel mass are both high
    (early-race conditions where the car is heavy and tires are fresh), the
    combined effect favors podium outcomes. This is physically consistent:
    top teams manage tire degradation more effectively under heavy fuel loads,
    maintaining lap time while competitors lose performance.
\end{enumerate}

\subsection{Machine Learning Validation (LOGO-CV)}

To test whether the statistical equation generalizes to unseen circuits, a
multinomial logistic regression with balanced class weighting was validated
using leave-one-group-out cross-validation across the three races.

\subsubsection{Classification Performance}

\begin{table}[htbp]
  \centering
  \caption{LOGO-CV Classification Report}
  \label{tab:classification}
  \begin{tabular}{lcccc}
    \toprule
    \textbf{Tier} & \textbf{Precision} & \textbf{Recall} &
    \textbf{F1-Score} & \textbf{Support (laps)} \\
    \midrule
    No Points     & 0.74 & 0.67 & 0.70 & 1,468 \\
    Podium        & 0.57 & 0.89 & 0.69 & 429 \\
    Point Scorer  & 0.49 & 0.45 & 0.47 & 983 \\
    \midrule
    \multicolumn{4}{l}{\textbf{Overall Accuracy}} & \textbf{62.36\%} \\
    \bottomrule
  \end{tabular}
\end{table}

\subsubsection{Interpretation}

\textbf{Podium recall (0.89).} The model correctly identified 89\% of all
actual podium-finishing laps across held-out circuits. This is the strongest
result and demonstrates that the combination of grid position control and
telemetry-based features captures top-tier performance signatures reliably.

\textbf{Podium precision (0.57).} When the model predicts a podium lap, it
is correct 57\% of the time. While this leaves room for improvement, it
significantly outperforms random guessing for a three-class problem
(chance $\approx$ 33\%).

\textbf{Point Scorer recall (0.45).} The model caught 45\% of
point-scoring laps. The F1 midfield is genuinely harder to model---DRS
trains, traffic, undercut strategies, and intra-team orders create variance
that telemetry alone cannot fully explain. Nevertheless, 45\% recall is well
above chance and indicates that the model captures a meaningful portion of
mid-field dynamics.

\textbf{Overall accuracy (62.36\%).} For a three-class problem with natural
class imbalance and cross-circuit validation (no data leakage), this
represents a credible result that confirms the statistical model generalizes
beyond its training data.
